% ==============================================================================
% MÉMOIRE DE FIN D'ÉTUDES — FICHIER UNIQUE COMPLET
% Titre  : Système d'Analyse de Sentiments des Commentaires Étudiants
% Auteur : Hissein Nasser
% Année  : 2025-2026
% Compilation : pdflatex (2 passes min) ou latexmk -pdf memoire_complet.tex
% ==============================================================================
\documentclass[12pt,a4paper,oneside]{report}

% ─────────────────────────────────────────────────────────────────────────────
% PACKAGES
% ─────────────────────────────────────────────────────────────────────────────
\usepackage[utf8]{inputenc}
\usepackage[T1]{fontenc}
\usepackage[french]{babel}
\usepackage{lmodern}
\usepackage[top=2.5cm,bottom=2.5cm,left=3cm,right=2.5cm]{geometry}
\usepackage{setspace}
\onehalfspacing

% Couleurs
\usepackage{xcolor}
\definecolor{bleuprincipal}{RGB}{31,73,125}
\definecolor{bleuclaire}{RGB}{46,134,193}
\definecolor{vertprincipal}{RGB}{31,138,112}
\definecolor{rougeprincipal}{RGB}{231,76,60}
\definecolor{griscode}{RGB}{248,249,250}
\definecolor{grisbordure}{RGB}{204,204,204}
\definecolor{orangeaccent}{RGB}{211,84,0}
\definecolor{jauneneutre}{RGB}{241,196,15}

% Graphiques
\usepackage{graphicx}
\usepackage{float}
\usepackage{caption}
\usepackage{subcaption}
\captionsetup{font=small,labelfont={bf,color=bleuprincipal}}

% Tableaux
\usepackage{booktabs}
\usepackage{tabularx}
\usepackage{multirow}
\usepackage{array}
\usepackage{longtable}
\usepackage{colortbl}

% TikZ
\usepackage{tikz}
\usepackage{pgfplots}
\pgfplotsset{compat=1.18}
\usetikzlibrary{shapes.geometric,arrows.meta,positioning,calc,
                shadows,fit,backgrounds,matrix,patterns}

% Listings (code source)
\usepackage{listings}
\usepackage{lstautogobble}

\lstdefinestyle{stylePython}{
  language=Python,
  backgroundcolor=\color{griscode},
  basicstyle=\ttfamily\footnotesize,
  keywordstyle=\color{bleuprincipal}\bfseries,
  stringstyle=\color{rougeprincipal},
  commentstyle=\color{vertprincipal}\itshape,
  numberstyle=\tiny\color{gray},
  numbers=left, numbersep=8pt,
  breaklines=true, frame=single,
  rulecolor=\color{grisbordure},
  tabsize=4, showstringspaces=false,
  captionpos=b, autogobble=true,
  xleftmargin=1cm,
  morekeywords={self,True,False,None,async,await,yield}
}

\lstdefinestyle{styleDart}{
  language=Java,
  backgroundcolor=\color{griscode},
  basicstyle=\ttfamily\footnotesize,
  keywordstyle=\color{bleuprincipal}\bfseries,
  stringstyle=\color{rougeprincipal},
  commentstyle=\color{vertprincipal}\itshape,
  numberstyle=\tiny\color{gray},
  numbers=left, numbersep=8pt,
  breaklines=true, frame=single,
  rulecolor=\color{grisbordure},
  tabsize=2, showstringspaces=false,
  captionpos=b, autogobble=true,
  xleftmargin=1cm,
  morekeywords={final,const,var,String,int,double,bool,void,
                Widget,StatelessWidget,StatefulWidget,BuildContext,
                Future,async,await,override,required,class}
}

\lstdefinestyle{styleSQL}{
  language=SQL,
  backgroundcolor=\color{griscode},
  basicstyle=\ttfamily\footnotesize,
  keywordstyle=\color{bleuprincipal}\bfseries,
  stringstyle=\color{rougeprincipal},
  commentstyle=\color{gray}\itshape,
  numberstyle=\tiny\color{gray},
  numbers=left, numbersep=8pt,
  breaklines=true, frame=single,
  rulecolor=\color{grisbordure},
  captionpos=b, xleftmargin=1cm
}

\lstdefinestyle{styleJSON}{
  language=Python,
  backgroundcolor=\color{griscode},
  basicstyle=\ttfamily\footnotesize,
  keywordstyle=\color{bleuprincipal},
  stringstyle=\color{vertprincipal},
  numberstyle=\tiny\color{gray},
  numbers=left, numbersep=8pt,
  breaklines=true, frame=single,
  rulecolor=\color{grisbordure},
  captionpos=b, xleftmargin=1cm
}

\lstset{style=stylePython}

% En-têtes / pieds de page
\usepackage{fancyhdr}
\pagestyle{fancy}
\fancyhf{}
\fancyhead[L]{\small\textcolor{bleuprincipal}{\leftmark}}
\fancyhead[R]{\small\textcolor{gray}{Analyse de Sentiments Étudiants}}
\fancyfoot[C]{\textcolor{gray}{\thepage}}
\fancyfoot[R]{\small\textcolor{gray}{Hissein Nasser — PFE 2026}}
\renewcommand{\headrulewidth}{0.5pt}
\renewcommand{\headrule}{\hbox to\headwidth{%
  \color{bleuprincipal}\leaders\hrule height\headrulewidth\hfill}}

% Style des chapitres
\usepackage{titlesec}
\titleformat{\chapter}[block]
  {\normalfont\huge\bfseries\color{bleuprincipal}}
  {\colorbox{bleuprincipal}{\textcolor{white}{\quad\thechapter\quad}}}
  {1em}{}[\vspace{-0.5em}\rule{\textwidth}{2pt}]
\titleformat{\section}
  {\normalfont\Large\bfseries\color{bleuprincipal}}{\thesection}{0.8em}{}
\titleformat{\subsection}
  {\normalfont\large\bfseries\color{bleuclaire}}{\thesubsection}{0.8em}{}
\titleformat{\subsubsection}
  {\normalfont\normalsize\bfseries\color{orangeaccent}}{\thesubsubsection}{0.8em}{}

% Hyperliens
\usepackage[
  colorlinks=true,
  linkcolor=bleuprincipal,
  citecolor=vertprincipal,
  urlcolor=bleuclaire,
  pdfauthor={Hissein Nasser},
  pdftitle={Systeme d'Analyse de Sentiments des Commentaires Etudiants},
  pdfsubject={Memoire de Fin d'Etudes}
]{hyperref}

% Bibliographie
\usepackage[backend=bibtex,style=ieee,sorting=none]{biblatex}

% Acronymes
\usepackage[acronym,toc]{glossaries}
\makeglossaries

% Utilitaires
\usepackage{amsmath,amssymb}
\usepackage{enumitem}
\usepackage{pifont}
\usepackage{tcolorbox}
\tcbuselibrary{skins,breakable}

% Boîtes personnalisées
\newtcolorbox{infobox}[1][]{enhanced,breakable,
  colback=bleuclaire!8!white,colframe=bleuclaire,
  fonttitle=\bfseries,title={#1},
  left=4pt,right=4pt,top=4pt,bottom=4pt}
\newtcolorbox{definitionbox}[1][]{enhanced,breakable,
  colback=vertprincipal!8!white,colframe=vertprincipal,
  fonttitle=\bfseries,title={#1},
  left=4pt,right=4pt,top=4pt,bottom=4pt}

% Commandes
\newcommand{\code}[1]{\texttt{\textcolor{bleuprincipal}{#1}}}

% ─────────────────────────────────────────────────────────────────────────────
% ACRONYMES
% ─────────────────────────────────────────────────────────────────────────────
\newacronym{nlp}{NLP}{Natural Language Processing}
\newacronym{taln}{TALN}{Traitement Automatique du Langage Naturel}
\newacronym{bert}{BERT}{Bidirectional Encoder Representations from Transformers}
\newacronym{api}{API}{Application Programming Interface}
\newacronym{rest}{REST}{Representational State Transfer}
\newacronym{orm}{ORM}{Object-Relational Mapping}
\newacronym{sgbd}{SGBD}{Système de Gestion de Bases de Données}
\newacronym{ml}{ML}{Machine Learning}
\newacronym{dl}{DL}{Deep Learning}
\newacronym{ia}{IA}{Intelligence Artificielle}
\newacronym{http}{HTTP}{HyperText Transfer Protocol}
\newacronym{json}{JSON}{JavaScript Object Notation}
\newacronym{sql}{SQL}{Structured Query Language}
\newacronym{cors}{CORS}{Cross-Origin Resource Sharing}
\newacronym{gpu}{GPU}{Graphics Processing Unit}
\newacronym{cpu}{CPU}{Central Processing Unit}
\newacronym{pfe}{PFE}{Projet de Fin d'Études}
\newacronym{uml}{UML}{Unified Modeling Language}
\newacronym{crud}{CRUD}{Create, Read, Update, Delete}
\newacronym{sdk}{SDK}{Software Development Kit}
\newacronym{mcd}{MCD}{Modèle Conceptuel de Données}
\newacronym{mld}{MLD}{Modèle Logique de Données}
\newacronym{ui}{UI}{User Interface}
\newacronym{bpe}{BPE}{Byte-Pair Encoding}

% ─────────────────────────────────────────────────────────────────────────────
% BIBLIOGRAPHIE INLINE
% ─────────────────────────────────────────────────────────────────────────────
\usepackage{filecontents}
\begin{filecontents}{refs.bib}
@article{martin2020camembert,
  author  = {Martin, Louis and others},
  title   = {{CamemBERT: a Tasty French Language Model}},
  journal = {Proceedings of ACL 2020},
  year    = {2020},
  pages   = {7203--7219}
}
@article{devlin2019bert,
  author  = {Devlin, Jacob and Chang, Ming-Wei and Lee, Kenton and Toutanova, Kristina},
  title   = {{BERT: Pre-training of Deep Bidirectional Transformers}},
  journal = {Proceedings of NAACL-HLT 2019},
  year    = {2019},
  pages   = {4171--4186}
}
@article{vaswani2017attention,
  author  = {Vaswani, Ashish and others},
  title   = {{Attention Is All You Need}},
  journal = {Advances in Neural Information Processing Systems},
  year    = {2017},
  volume  = {30}
}
@article{liu2019roberta,
  author  = {Liu, Yinhan and others},
  title   = {{RoBERTa: A Robustly Optimized BERT Pretraining Approach}},
  journal = {arXiv:1907.11692},
  year    = {2019}
}
@inproceedings{kudo2018sentencepiece,
  author    = {Kudo, Taku and Richardson, John},
  title     = {{SentencePiece: A simple and language independent subword tokenizer}},
  booktitle = {EMNLP 2018 System Demonstrations},
  year      = {2018},
  pages     = {66--71}
}
@misc{fastapi2024,
  author       = {Ramirez, Sebastian},
  title        = {{FastAPI Framework}},
  year         = {2024},
  howpublished = {\url{https://fastapi.tiangolo.com}}
}
@misc{pytorch2024,
  author       = {{PyTorch Team}},
  title        = {{PyTorch Deep Learning Library}},
  year         = {2024},
  howpublished = {\url{https://pytorch.org}}
}
@misc{flutter2024,
  author       = {{Google LLC}},
  title        = {{Flutter SDK}},
  year         = {2024},
  howpublished = {\url{https://flutter.dev}}
}
@article{loshchilov2019adamw,
  author  = {Loshchilov, Ilya and Hutter, Frank},
  title   = {{Decoupled Weight Decay Regularization}},
  journal = {ICLR 2019},
  year    = {2019}
}
@article{pang2008opinion,
  author  = {Pang, Bo and Lee, Lillian},
  title   = {{Opinion Mining and Sentiment Analysis}},
  journal = {Foundations and Trends in Information Retrieval},
  year    = {2008},
  volume  = {2},
  pages   = {1--135}
}
\end{filecontents}
\addbibresource{refs.bib}

% ==============================================================================
%                          DÉBUT DU DOCUMENT
% ==============================================================================
\begin{document}
\pagenumbering{gobble}

% ══════════════════════════════════════════════════════════════════════════════
%  PAGE DE GARDE
% ══════════════════════════════════════════════════════════════════════════════
\begin{titlepage}
\thispagestyle{empty}
\begin{tikzpicture}[remember picture,overlay]
  \fill[bleuprincipal]   (current page.north west) rectangle ([yshift=-3.5cm]current page.north east);
  \fill[bleuprincipal!80!black] (current page.south west) rectangle ([yshift=1.8cm]current page.south east);
  \fill[vertprincipal]   (current page.north west) rectangle ([xshift=0.8cm]current page.south west);
  \fill[jauneneutre]     ([xshift=0.8cm]current page.north west) rectangle ([xshift=1.1cm]current page.south west);
\end{tikzpicture}

\vspace*{0.5cm}
\begin{center}
{\color{white}\Large\bfseries UNIVERSITÉ --- DÉPARTEMENT INFORMATIQUE}\\[0.3cm]
{\color{white!80!gray}\large Filière : Génie Informatique}\\[0.2cm]
{\color{white!80!gray}\normalsize Option : Intelligence Artificielle \& Génie Logiciel}

\vspace{1.2cm}
\hrule height 1pt\relax
\vspace{0.3cm}
{\color{bleuprincipal!30!black}\small MÉMOIRE DE FIN D'ÉTUDES}\\
{\color{bleuprincipal!30!black}\small Pour l'obtention du diplôme de Licence / Master en Informatique}
\vspace{0.3cm}
\hrule height 1pt\relax

\vspace{1.5cm}
\begin{tcolorbox}[enhanced,colback=bleuprincipal!10!white,colframe=bleuprincipal,
  boxrule=2pt,arc=8pt,left=12pt,right=12pt,top=10pt,bottom=10pt,width=14cm]
\begin{center}
{\huge\bfseries\color{bleuprincipal}
Système d'Analyse de\\[0.3em]
Sentiments des\\[0.3em]
Commentaires Étudiants\\[0.3em]
en Temps Réel}\\[0.8em]
{\large\color{gray}Basé sur CamemBERT, FastAPI, PostgreSQL et Flutter}
\end{center}
\end{tcolorbox}

\vspace{1.8cm}
\begin{minipage}[t]{0.45\textwidth}
\begin{flushleft}
\textbf{\color{bleuprincipal}Réalisé par :}\\[0.5em]
\large Hissein Nasser\\
\small Étudiant en Informatique
\end{flushleft}
\end{minipage}
\hfill
\begin{minipage}[t]{0.45\textwidth}
\begin{flushright}
\textbf{\color{bleuprincipal}Encadré par :}\\[0.5em]
\large M./Mme. [Nom Encadreur]\\
\small Département Informatique
\end{flushright}
\end{minipage}

\vspace{1.8cm}
\begin{tikzpicture}
  \node[fill=bleuprincipal!10,rounded corners=5pt,inner sep=6pt] at (0,0)   {\small\color{bleuprincipal}\textbf{CamemBERT}};
  \node[fill=vertprincipal!15,rounded corners=5pt,inner sep=6pt] at (3,0)   {\small\color{vertprincipal}\textbf{FastAPI}};
  \node[fill=rougeprincipal!10,rounded corners=5pt,inner sep=6pt] at (6,0)  {\small\color{rougeprincipal}\textbf{PostgreSQL}};
  \node[fill=bleuclaire!15,rounded corners=5pt,inner sep=6pt] at (9,0)      {\small\color{bleuclaire}\textbf{Flutter}};
  \node[fill=orangeaccent!10,rounded corners=5pt,inner sep=6pt] at (1.5,-0.9){\small\color{orangeaccent}\textbf{PyTorch}};
  \node[fill=bleuprincipal!10,rounded corners=5pt,inner sep=6pt] at (4.5,-0.9){\small\color{bleuprincipal}\textbf{Streamlit}};
  \node[fill=vertprincipal!15,rounded corners=5pt,inner sep=6pt] at (7.5,-0.9){\small\color{vertprincipal}\textbf{SQLAlchemy}};
\end{tikzpicture}

\vspace{1.5cm}
{\color{white}\large\bfseries Année Académique 2025 -- 2026}
\end{center}
\end{titlepage}
\cleardoublepage

% ══════════════════════════════════════════════════════════════════════════════
%  PAGES LIMINAIRES
% ══════════════════════════════════════════════════════════════════════════════
\pagenumbering{roman}\setcounter{page}{1}

% ── Dédicace ──────────────────────────────────────────────────────────────────
\chapter*{Dédicace}
\addcontentsline{toc}{chapter}{Dédicace}
\thispagestyle{empty}
\vspace*{3cm}
\begin{flushright}
\begin{minipage}{0.7\textwidth}
\begin{center}
{\Large\color{bleuprincipal}\textit{\og À ma famille, \fg{}}}\\[1.5cm]
\normalsize\itshape
À mes parents, pour leur amour inconditionnel, leurs sacrifices
infinis et leur confiance inébranlable en moi.\\[1cm]
À mes frères et sœurs, compagnons de vie, dont le soutien
moral et l'encouragement ont rendu ce travail possible.\\[1cm]
À tous ceux qui croient que la technologie, au service
de l'éducation, peut transformer le monde.\\[2cm]
{\large\bfseries\color{bleuprincipal}--- Hissein Nasser ---}
\end{center}
\end{minipage}
\end{flushright}
\cleardoublepage

% ── Remerciements ─────────────────────────────────────────────────────────────
\chapter*{Remerciements}
\addcontentsline{toc}{chapter}{Remerciements}

La réalisation de ce projet de fin d'études n'aurait pas été possible sans
le concours de nombreuses personnes à qui je tiens à exprimer ma sincère gratitude.

\vspace{0.5cm}
En premier lieu, mes remerciements les plus chaleureux vont à mon encadreur,
\textbf{M./Mme. [Nom de l'Encadreur]}, pour son suivi rigoureux, sa
disponibilité constante et ses précieux conseils scientifiques.

\vspace{0.5cm}
Je tiens également à remercier les membres du jury qui ont accepté d'évaluer
ce mémoire et de consacrer leur temps à l'examen de ce travail.

\vspace{0.5cm}
Ma gratitude va aussi à l'ensemble du corps enseignant du département
Informatique, dont l'enseignement de qualité tout au long de ma formation
m'a fourni les bases théoriques et pratiques indispensables.

\vspace{0.5cm}
Je remercie la communauté open-source, en particulier les équipes d'Inria
et CNRS pour CamemBERT, l'équipe HuggingFace pour Transformers, ainsi que
toutes les communautés qui contribuent à l'avancement de l'intelligence artificielle.

\vspace{0.5cm}
Enfin, un remerciement profond et sincère à ma famille pour leur patience
et leur soutien indéfectible tout au long de ces années d'études.

\begin{flushright}\textit{Hissein Nasser --- Mai 2026}\end{flushright}
\cleardoublepage

% ── Résumé ────────────────────────────────────────────────────────────────────
\chapter*{Résumé}
\addcontentsline{toc}{chapter}{Résumé}
\begin{tcolorbox}[enhanced,colback=bleuprincipal!5!white,colframe=bleuprincipal,
  boxrule=1.5pt,arc=6pt,left=10pt,right=10pt,top=8pt,bottom=8pt]
Ce mémoire présente la conception et le développement d'un système complet
d'analyse automatique de sentiments appliqué aux commentaires d'étudiants
sur leurs cours universitaires. Le système repose sur le modèle de langage
\textbf{CamemBERT} (\textit{camembert-base}), fine-tuné pour la classification
de sentiments en trois catégories~: \textbf{Négatif}, \textbf{Neutre} et
\textbf{Positif}, sur un jeu de données de 225 commentaires étudiants en
français avec les hyperparamètres~: taux d'apprentissage $2\times10^{-5}$,
batch 16, 3 epochs, 128 tokens maximum. Le modèle atteint une précision de
\textbf{79,41\%} sur le jeu de test. L'architecture globale comprend~:
un backend RESTful FastAPI avec PostgreSQL, un tableau de bord Streamlit
et une application mobile Flutter.
\end{tcolorbox}
\vspace{0.5em}
\noindent\textbf{Mots-clés~:} Analyse de sentiments, TALN, CamemBERT,
Apprentissage automatique, FastAPI, Flutter, PostgreSQL, Fine-tuning.
\cleardoublepage

% ── Abstract ──────────────────────────────────────────────────────────────────
\chapter*{Abstract}
\addcontentsline{toc}{chapter}{Abstract}
\begin{tcolorbox}[enhanced,colback=vertprincipal!5!white,colframe=vertprincipal,
  boxrule=1.5pt,arc=6pt,left=10pt,right=10pt,top=8pt,bottom=8pt]
This thesis presents the design and development of a complete automated
sentiment analysis system applied to student feedback on university courses.
The system is built upon the \textbf{CamemBERT} language model fine-tuned
for three-class sentiment classification: \textbf{Negative}, \textbf{Neutral},
and \textbf{Positive}, trained on 225 French student comments with
hyperparameters: learning rate $2\times10^{-5}$, batch 16, 3 epochs,
max 128 tokens. The model achieves \textbf{79.41\%} accuracy on the test set.
The architecture includes a FastAPI REST backend with PostgreSQL, a Streamlit
analytics dashboard, and a Flutter mobile application.
\end{tcolorbox}
\vspace{0.5em}
\noindent\textbf{Keywords~:} Sentiment analysis, NLP, CamemBERT, Machine
learning, FastAPI, Flutter, PostgreSQL, Fine-tuning, Transfer learning.
\cleardoublepage

% ── Tables ────────────────────────────────────────────────────────────────────
{\hypersetup{linkcolor=black}
\tableofcontents\cleardoublepage
\listoffigures\cleardoublepage
\listoftables\cleardoublepage}
\printglossary[type=\acronymtype,title={Liste des Acronymes}]
\cleardoublepage

% ══════════════════════════════════════════════════════════════════════════════
%  CORPS DU RAPPORT
% ══════════════════════════════════════════════════════════════════════════════
\pagenumbering{arabic}\setcounter{page}{1}

% ══════════════════════════════════════════════════════════════════════════════
%  INTRODUCTION GÉNÉRALE
% ══════════════════════════════════════════════════════════════════════════════
\chapter*{Introduction Générale}
\addcontentsline{toc}{chapter}{Introduction Générale}
\markboth{INTRODUCTION GÉNÉRALE}{}

\begin{infobox}[Contexte de l'étude]
L'essor du numérique dans l'enseignement supérieur génère chaque année
des millions de retours étudiants. Traiter ces données manuellement est
devenu impossible. L'intelligence artificielle offre une réponse concrète
et immédiate.
\end{infobox}

\section*{Contexte général}

L'enseignement supérieur traverse une transformation profonde depuis
l'avènement des technologies numériques. La massification de l'éducation
universitaire, conjuguée à l'explosion des plateformes d'apprentissage en
ligne, a multiplié exponentiellement les interactions numériques entre
enseignants et étudiants. Selon plusieurs études pédagogiques récentes,
un étudiant génère en moyenne plusieurs dizaines de commentaires par
semestre. À l'échelle d'un département, cela représente des milliers de
textes à analyser chaque année.

\section*{Problématique}

Comment peut-on concevoir un système robuste et précis permettant d'analyser
automatiquement et en temps réel les sentiments exprimés dans les commentaires
d'étudiants en français, et de mettre ces analyses à la disposition des
enseignants de manière exploitable ?

\section*{Solution proposée}

Ce projet propose une réponse complète à travers un système multi-couches
intégrant~: un modèle d'\acrfull{ia} basé sur \textbf{CamemBERT}, un
serveur \acrfull{api} RESTful FastAPI, un tableau de bord analytique Streamlit
et une application mobile Android Flutter.

\section*{Objectifs du projet}

\begin{enumerate}[leftmargin=2.5em]
  \item Constituer un jeu de données de commentaires étudiants en français,
        équilibré entre les classes de sentiment.
  \item Fine-tuner CamemBERT pour la classification ternaire de sentiments.
  \item Développer une API REST robuste et documentée.
  \item Concevoir un tableau de bord analytique interactif.
  \item Développer une application mobile multiplateforme.
  \item Évaluer les performances et identifier les pistes d'amélioration.
\end{enumerate}

\section*{Organisation du mémoire}

\noindent\textbf{Chapitre 1} --- Contexte et Problématique : contexte
général, limites des approches existantes et justification technologique.\\
\noindent\textbf{Chapitre 2} --- Analyse et Conception : diagrammes UML,
architecture globale, modèle de données.\\
\noindent\textbf{Chapitre 3} --- Implémentation : environnements, code source
commenté, algorithmes clés.\\
\noindent\textbf{Chapitre 4} --- Tests et Résultats : métriques de performance,
scénarios de test, analyse.\\
\noindent\textbf{Chapitre 5} --- Discussion et Perspectives : analyse critique,
limites et évolutions futures.
\cleardoublepage

% ══════════════════════════════════════════════════════════════════════════════
%  CHAPITRE 1 : CONTEXTE ET PROBLÉMATIQUE
% ══════════════════════════════════════════════════════════════════════════════
\chapter{Contexte et Problématique}

\begin{infobox}[Objectif du chapitre]
Ce chapitre établit le cadre général de l'étude~: le contexte éducatif
numérique, la problématique du traitement des retours étudiants, l'étude
des solutions existantes et la valeur ajoutée de notre approche.
\end{infobox}

\section{Contexte général : La Révolution Numérique dans l'Éducation}

L'enseignement supérieur traverse une transformation profonde depuis l'avènement
des technologies numériques. La massification de l'éducation universitaire et
l'explosion des plateformes d'apprentissage en ligne ont multiplié les
interactions numériques entre enseignants et étudiants.

\subsection{L'importance du retour étudiant}

Les retours des étudiants constituent un indicateur fondamental de la qualité
pédagogique. Ils permettent aux enseignants d'identifier les points faibles
de leurs cours, d'adapter leurs méthodes et d'améliorer l'expérience
d'apprentissage. Cependant, leur valeur réelle ne peut être exploitée que si
ces données sont analysées de manière systématique, objective et rapide.

\subsection{Limitations du traitement manuel}

\begin{table}[H]
\centering
\caption{Limitations du traitement manuel des commentaires}
\label{tab:limites_manuel}
\renewcommand{\arraystretch}{1.4}
\begin{tabular}{>{\bfseries}p{3.2cm} p{9cm}}
\toprule
\rowcolor{bleuprincipal!15}
\textbf{Limitation} & \textbf{Impact} \\
\midrule
Lenteur & Délai de plusieurs semaines, insights obsolètes \\
\rowcolor{gray!8}
Subjectivité & Biais cognitifs, interprétation variable \\
Coût élevé & Mobilisation de ressources humaines importantes \\
\rowcolor{gray!8}
Non scalable & Impossible avec l'augmentation des effectifs \\
Exhaustivité & Tendance à ignorer les nuances \\
\bottomrule
\end{tabular}
\end{table}

\section{Problématique Détaillée}

\subsection{Dimension linguistique}

Le français présente des défis spécifiques pour l'analyse automatique de
sentiments~: richesse morphologique, usage fréquent de la négation (``ce cours
n'est pas mauvais'' exprime quelque chose de positif), intensification
(``vraiment excellent'', ``assez décevant'') et ironie (``super, encore un cours
incompréhensible...'').

\subsection{Dimension technologique}

Les approches traditionnelles basées sur des dictionnaires de polarité ou des
modèles statistiques simples (SVM, Naive Bayes) présentent des performances
insuffisantes sur le français informel étudiant. Un modèle Transformer profond
pré-entraîné sur de larges corpus français s'impose naturellement.

\section{Étude de l'Existant}

\subsection{Comparaison des approches}

\begin{table}[H]
\centering
\caption{Comparaison des approches d'analyse de sentiments pour le français}
\label{tab:comparaison}
\renewcommand{\arraystretch}{1.4}
\begin{tabular}{lcccc}
\toprule
\rowcolor{bleuprincipal!15}
\textbf{Approche} & \textbf{Précision} & \textbf{Gratuit} & \textbf{Hors-ligne} & \textbf{Personnalisable} \\
\midrule
VADER (anglais) & $<$55\% & Oui & Oui & Limitée \\
\rowcolor{gray!8}
Google NL API & $\sim$72\% & Non & Non & Non \\
mBERT & $\sim$70\% & Oui & Oui & Oui \\
\rowcolor{gray!8}
XLM-RoBERTa & $\sim$74\% & Oui & Oui & Oui \\
\rowcolor{vertprincipal!15}
\textbf{CamemBERT (notre)} & \textbf{79,41\%} & \textbf{Oui} & \textbf{Oui} & \textbf{Oui} \\
\bottomrule
\end{tabular}
\end{table}

\section{Fondements Théoriques}

\subsection{L'architecture Transformer}

Le Transformer \cite{vaswani2017attention} introduit le mécanisme
d'\textbf{attention multi-têtes} permettant au modèle de pondérer l'importance
de chaque mot dans une séquence par rapport à tous les autres, capturant
des dépendances longue-distance inaccessibles aux RNN/LSTM.

\begin{equation}
\text{Attention}(Q,K,V) = \text{softmax}\!\left(\frac{QK^T}{\sqrt{d_k}}\right)V
\label{eq:attention}
\end{equation}

\subsection{BERT et CamemBERT}

\acrfull{bert} \cite{devlin2019bert} exploite un contexte bidirectionnel via
deux tâches de pré-entraînement~: Masked Language Modeling (15\% des tokens
masqués) et Next Sentence Prediction. CamemBERT \cite{martin2020camembert}
reprend l'architecture RoBERTa \cite{liu2019roberta} et est entraîné
uniquement sur du texte français (138 Go, corpus OSCAR). Il utilise la
tokenisation SentencePiece \cite{kudo2018sentencepiece} avec un vocabulaire
de 32\,000 tokens.

\subsection{Le fine-tuning par transfert d'apprentissage}

Pour la classification, une couche linéaire est ajoutée sur le token
\texttt{[CLS]} de CamemBERT~:
\begin{equation}
P(\text{classe}_k|\mathbf{x}) = \text{softmax}(W\cdot\mathbf{h}_{CLS}+b)_k
\label{eq:classification}
\end{equation}

\section{Valeur Ajoutée du Projet}

\begin{definitionbox}[Contributions principales]
\begin{enumerate}
  \item \textbf{Spécialisation domaine~:} fine-tuning sur des commentaires
        académiques français.
  \item \textbf{Système complet bout en bout~:} de la collecte à l'affichage
        mobile.
  \item \textbf{Temps réel~:} analyse et retour en moins d'une seconde.
  \item \textbf{Accessibilité~:} application mobile pour les étudiants,
        dashboard pour les enseignants.
  \item \textbf{Gratuité et confidentialité~:} déployable entièrement en local.
\end{enumerate}
\end{definitionbox}

\section*{Conclusion du Chapitre 1}
Ce chapitre a établi le contexte motivant notre projet, analysé les limites
des approches existantes et positionné CamemBERT comme le choix technologique
le plus pertinent. Le chapitre suivant détaille la conception architecturale.
\cleardoublepage

% ══════════════════════════════════════════════════════════════════════════════
%  CHAPITRE 2 : ANALYSE ET CONCEPTION
% ══════════════════════════════════════════════════════════════════════════════
\chapter{Analyse et Conception}

\begin{infobox}[Objectif du chapitre]
Analyse fonctionnelle et non fonctionnelle, diagrammes UML, architecture
globale, modèle de données et justification des choix technologiques.
\end{infobox}

\section{Analyse des Besoins}

\subsection{Acteurs du système}

\begin{table}[H]
\centering
\caption{Acteurs du système}
\label{tab:acteurs}
\renewcommand{\arraystretch}{1.4}
\begin{tabular}{>{\bfseries}p{2.8cm} p{4.5cm} p{5.5cm}}
\toprule
\rowcolor{bleuprincipal!15}
\textbf{Acteur} & \textbf{Rôle} & \textbf{Interactions principales} \\
\midrule
Étudiant & Soumetteur de commentaires & Soumettre, consulter résultat, voir historique \\
\rowcolor{gray!8}
Enseignant & Consommateur d'analyses & Dashboard, filtres, statistiques \\
Système IA & Acteur interne & Analyser, retourner les scores \\
\bottomrule
\end{tabular}
\end{table}

\subsection{Besoins fonctionnels}

\begin{enumerate}[label=\textbf{BF\arabic*.},leftmargin=3em]
  \item Soumettre un commentaire textuel d'au moins 5 caractères.
  \item Analyser automatiquement le sentiment et retourner la classe + scores.
  \item Associer un commentaire à une matière académique.
  \item Sauvegarder chaque analyse dans la base de données.
  \item Consulter l'historique paginé des commentaires.
  \item Filtrer par sentiment, matière et date.
  \item Supprimer un commentaire par son identifiant.
  \item Exposer des statistiques agrégées.
  \item Navigation mobile entre écran d'analyse et historique.
  \item Visualisations interactives dans le dashboard.
\end{enumerate}

\subsection{Besoins non fonctionnels}

\begin{table}[H]
\centering
\caption{Besoins non fonctionnels}
\label{tab:bnf}
\renewcommand{\arraystretch}{1.4}
\begin{tabular}{>{\bfseries}p{2.8cm} p{10cm}}
\toprule
\rowcolor{bleuprincipal!15}
\textbf{Critère} & \textbf{Exigence} \\
\midrule
Performance & Temps de réponse d'une analyse $<$ 2 secondes \\
\rowcolor{gray!8}
Scalabilité & Ajout de matières sans modification du code \\
Maintenabilité & Code modulaire, documenté, versionné via Git \\
\rowcolor{gray!8}
Sécurité & Données stockées localement, sans transfert cloud \\
Compatibilité & Application mobile Android 8.0+ \\
\rowcolor{gray!8}
Portabilité & Déployable sous Ubuntu sans dépendances propriétaires \\
\bottomrule
\end{tabular}
\end{table}

\section{Diagrammes UML}

\subsection{Diagramme de cas d'utilisation}

\begin{figure}[H]
\centering
\begin{tikzpicture}[
  actor/.style={circle,draw=bleuprincipal,fill=bleuprincipal!10,minimum size=1cm,font=\small\bfseries},
  usecase/.style={ellipse,draw=bleuprincipal,fill=white,minimum width=3.8cm,minimum height=0.9cm,font=\footnotesize,align=center,thick}]
  \node[actor] (E) at (-4.5,0) {};
  \node[below=0.1cm of E,font=\small\bfseries] {Étudiant};
  \node[actor] (P) at (7,0) {};
  \node[below=0.1cm of P,font=\small\bfseries] {Enseignant};
  \draw[dashed,bleuprincipal!50,thick,rounded corners=8pt] (-2,-4) rectangle (5,4);
  \node[font=\small\bfseries\color{bleuprincipal}] at (1.5,3.6) {Système};
  \node[usecase] (u1) at (1.5, 3.0) {Soumettre un\\commentaire};
  \node[usecase] (u2) at (1.5, 1.8) {Voir résultat\\d'analyse};
  \node[usecase] (u3) at (1.5, 0.6) {Consulter\\l'historique};
  \node[usecase] (u4) at (1.5,-0.6) {Filtrer les\\commentaires};
  \node[usecase] (u5) at (1.5,-1.8) {Supprimer un\\commentaire};
  \node[usecase] (u6) at (1.5,-3.0) {Voir les\\statistiques};
  \foreach \u in {u1,u2,u3,u4,u5} \draw[->] (E)--(\u);
  \foreach \u in {u3,u4,u5,u6} \draw[->] (P)--(\u);
\end{tikzpicture}
\caption{Diagramme de cas d'utilisation}
\label{fig:usecase}
\end{figure}

\subsection{Diagramme de séquence — Analyse d'un commentaire}

\begin{figure}[H]
\centering
\begin{tikzpicture}[
  box/.style={rectangle,draw=bleuprincipal,fill=bleuprincipal!10,minimum width=2.4cm,minimum height=0.7cm,font=\small\bfseries,align=center},
  msg/.style={->,thick,color=bleuprincipal!80},
  ret/.style={->,thick,dashed,color=vertprincipal}]
  \node[box] (A) at (0,0)    {Flutter};
  \node[box] (B) at (3.5,0)  {FastAPI};
  \node[box] (C) at (7,0)    {Service\\Prédiction};
  \node[box] (D) at (10.5,0) {CamemBERT};
  \node[box] (E) at (14,0)   {PostgreSQL};
  \foreach \x in {0,3.5,7,10.5,14} \draw[dashed,bleuprincipal!40] (\x,-0.35)--(\x,-9);
  \draw[msg] (0,-1.0)--node[above,font=\scriptsize]{POST /api/predict}(3.5,-1.0);
  \draw[msg] (3.5,-1.8)--node[above,font=\scriptsize]{predict(texte)}(7,-1.8);
  \draw[msg] (7,-2.6)--node[above,font=\scriptsize]{tokeniser()}(10.5,-2.6);
  \draw[msg] (10.5,-3.4)--node[above,font=\scriptsize]{forward()}(10.5,-3.4);
  \draw[ret] (10.5,-4.2)--node[above,font=\scriptsize]{logits}(7,-4.2);
  \draw[msg] (7,-5.0)--node[above,font=\scriptsize]{softmax()}(7,-5.0);
  \draw[ret] (7,-5.8)--node[above,font=\scriptsize]{\{sentiment,scores\}}(3.5,-5.8);
  \draw[msg] (3.5,-6.6)--node[above,font=\scriptsize]{INSERT}(14,-6.6);
  \draw[ret] (14,-7.4)--node[above,font=\scriptsize]{id,date}(3.5,-7.4);
  \draw[ret] (3.5,-8.2)--node[above,font=\scriptsize]{JSON résultat}(0,-8.2);
\end{tikzpicture}
\caption{Diagramme de séquence : analyse d'un commentaire}
\label{fig:sequence}
\end{figure}

\section{Architecture Globale du Système}

\begin{figure}[H]
\centering
\begin{tikzpicture}[
  layer/.style={rectangle,rounded corners=6pt,minimum width=3.5cm,minimum height=1.1cm,
                align=center,font=\small\bfseries,thick,draw},
  arr/.style={->,thick,color=bleuprincipal!70}]
  \node[layer,fill=bleuclaire!20,draw=bleuclaire]   (F) at (0,5)   {Application\\Flutter};
  \node[layer,fill=bleuclaire!20,draw=bleuclaire]   (S) at (5,5)   {Dashboard\\Streamlit};
  \node[layer,fill=bleuprincipal!15,draw=bleuprincipal,minimum width=9cm] (API) at (2.5,3.2) {API REST --- FastAPI (port 8000)};
  \node[layer,fill=vertprincipal!15,draw=vertprincipal] (SP) at (0,1.4) {Service\\Prédiction};
  \node[layer,fill=vertprincipal!15,draw=vertprincipal] (R)  at (5,1.4) {Routeur\\Commentaires};
  \node[layer,fill=orangeaccent!15,draw=orangeaccent]  (CB) at (0,-0.4) {CamemBERT\\Fine-tuné};
  \node[layer,fill=rougeprincipal!10,draw=rougeprincipal] (DB) at (5,-0.4) {PostgreSQL};
  \draw[arr] (F)--(API); \draw[arr] (S)--(API);
  \draw[arr] (API)--(SP); \draw[arr] (API)--(R);
  \draw[arr] (SP)--(CB); \draw[arr] (R)--(DB);
  \node[font=\tiny,gray] at (2.5,-1.4) {4 couches : Présentation --- API --- Métier --- Données};
\end{tikzpicture}
\caption{Architecture globale en 4 couches}
\label{fig:architecture}
\end{figure}

\section{Modèle de Données}

\subsection{Table principale — commentaires}

\begin{lstlisting}[style=styleSQL,caption={Schéma SQL de la table commentaires},label={lst:sql}]
CREATE TABLE commentaires (
    id               SERIAL           PRIMARY KEY,
    texte_original   TEXT             NOT NULL,
    texte_nettoye    TEXT,
    matiere          VARCHAR(100),
    sentiment_predit INTEGER          NOT NULL, -- 0, 1 ou 2
    label_sentiment  VARCHAR(20)      NOT NULL, -- 'Negatif','Neutre','Positif'
    score_negatif    DOUBLE PRECISION,
    score_neutre     DOUBLE PRECISION,
    score_positif    DOUBLE PRECISION,
    date_creation    TIMESTAMP        DEFAULT NOW()
);
\end{lstlisting}

\section{Justification des Choix Technologiques}

\begin{table}[H]
\centering
\caption{Justification des choix technologiques}
\label{tab:choix}
\renewcommand{\arraystretch}{1.5}
\begin{tabular}{p{2.4cm} p{2.8cm} p{7.5cm}}
\toprule
\rowcolor{bleuprincipal!15}
\textbf{Technologie} & \textbf{Rôle} & \textbf{Justification} \\
\midrule
CamemBERT & Modèle NLP & État de l'art pour le français, 138 Go d'entraînement \\
\rowcolor{gray!8}
PyTorch & Framework DL & Standard industriel, support HuggingFace natif \\
FastAPI & API REST & Plus rapide des frameworks Python, typage fort, Swagger auto \\
\rowcolor{gray!8}
SQLAlchemy & ORM & Abstraction DB, sécurité injection SQL \\
PostgreSQL & Base de données & ACID, performances, open-source \\
\rowcolor{gray!8}
Streamlit & Dashboard & Python pur, Plotly natif, déploiement instantané \\
Flutter & App mobile & Cross-platform Android/iOS, Material 3, performances natives \\
\bottomrule
\end{tabular}
\end{table}

\section*{Conclusion du Chapitre 2}
L'analyse fonctionnelle a permis d'identifier les besoins des utilisateurs
et de concevoir une architecture 4 couches répondant à ces besoins. Le
chapitre suivant détaille l'implémentation concrète.
\cleardoublepage

% ══════════════════════════════════════════════════════════════════════════════
%  CHAPITRE 3 : IMPLÉMENTATION
% ══════════════════════════════════════════════════════════════════════════════
\chapter{Implémentation}

\begin{infobox}[Objectif du chapitre]
Environnement de développement, structure du projet, algorithmes clés
et parties importantes du code source.
\end{infobox}

\section{Environnements de Développement}

\begin{table}[H]
\centering
\caption{Logiciels et versions utilisés}
\label{tab:software}
\renewcommand{\arraystretch}{1.4}
\begin{tabular}{lll}
\toprule
\rowcolor{bleuprincipal!15}
\textbf{Logiciel} & \textbf{Version} & \textbf{Rôle} \\
\midrule
Python & 3.12 & Langage backend et IA \\
\rowcolor{gray!8}
PyTorch & 2.x & Framework Deep Learning \\
Transformers (HuggingFace) & 4.x & Modèles pré-entraînés \\
\rowcolor{gray!8}
FastAPI & 0.110+ & Framework API REST \\
SQLAlchemy & 2.0+ & ORM Python-PostgreSQL \\
\rowcolor{gray!8}
PostgreSQL & 16 & Base de données relationnelle \\
Flutter SDK & 3.24.5 & Application mobile \\
\rowcolor{gray!8}
Android Studio & 2024.2 & Émulateur Android \\
Ubuntu & 24.04 LTS & Système d'exploitation \\
\bottomrule
\end{tabular}
\end{table}

\section{Structure du Projet}

\begin{lstlisting}[language=bash,caption={Structure complète du projet},
  basicstyle=\ttfamily\footnotesize,frame=single,
  backgroundcolor=\color{griscode},numbers=none,xleftmargin=0.5cm]
sentiment_etudiants/
|-- ai_model/
|   |-- preprocessing/
|   |   |-- nettoyage_texte.py   # Pipeline NLP de nettoyage
|   |   |-- tokenisation.py      # Tokenisation CamemBERT
|   |   `-- preparer_donnees.py  # Preparation du dataset
|   |-- modele_camembert.py      # Chargement/sauvegarde modele
|   |-- entrainement.py          # Fine-tuning (3 epochs)
|   |-- evaluation.py            # Accuracy, F1, matrice confusion
|   `-- saved_model/             # Modele fine-tune
|-- backend/
|   |-- database.py              # Connexion PostgreSQL / SQLAlchemy
|   |-- models/modeles_db.py     # ORM : table commentaires
|   |-- routers/commentaires.py  # 5 routes REST
|   |-- services/prediction_service.py  # Singleton inference
|   `-- main.py                  # Point d'entree FastAPI
|-- dashboard/
|   |-- app.py                   # Page accueil Streamlit
|   |-- components/api_client.py # Client HTTP
|   `-- pages/
|       |-- 1_Analyse.py
|       |-- 2_Statistiques.py
|       `-- 3_Historique.py
|-- mobile_app/lib/
|   |-- main.dart                # Navigation + MaterialApp
|   |-- models/commentaire.dart  # Modele de donnees Dart
|   |-- services/api_service.dart# Client HTTP Dart
|   |-- widgets/sentiment_card.dart # Widget reutilisable
|   `-- screens/
|       |-- analyse_screen.dart
|       `-- historique_screen.dart
`-- datasets/
    |-- raw/commentaires_etudiants.csv  # 225 commentaires
    `-- processed/{train,validation,test}.csv
\end{lstlisting}

\section{Module IA --- Pipeline CamemBERT}

\subsection{Pipeline de nettoyage NLP}

\begin{lstlisting}[style=stylePython,caption={Pipeline de nettoyage du texte},label={lst:nettoyage}]
def nettoyer_commentaire(texte: str) -> str:
    """Pipeline complet : normalisation, suppression bruit, minuscules."""
    if not isinstance(texte, str) or len(texte.strip()) == 0:
        return ""
    # 1. Normalisation Unicode NFC + guillemets typographiques
    texte = unicodedata.normalize("NFC", texte)
    texte = texte.replace(u"«", '"').replace(u"»", '"')
    # 2. Suppression URLs, emails, caracteres speciaux
    texte = re.sub(r"http\S+|www\S+", "", texte)
    texte = re.sub(r"\S+@\S+", "", texte)
    texte = re.sub(r"[^\w\sA-Za-zÀ-ž'',;:!?.\-]", " ", texte)
    # 3. Minuscules + normalisation espaces
    texte = texte.lower()
    texte = re.sub(r"\s+", " ", texte).strip()
    return texte
\end{lstlisting}

\subsection{Jeu de données}

\begin{table}[H]
\centering
\caption{Répartition et découpage du dataset}
\label{tab:dataset}
\renewcommand{\arraystretch}{1.4}
\begin{tabular}{lcccc}
\toprule
\rowcolor{bleuprincipal!15}
\textbf{Ensemble} & \textbf{Taille} & \textbf{Proportion} & \textbf{Négatif} & \textbf{Neutre / Positif} \\
\midrule
Entraînement & 157 & $\sim$70\% & 52 & 53 / 52 \\
\rowcolor{gray!8}
Validation & 34 & $\sim$15\% & 11 & 12 / 11 \\
Test & 34 & $\sim$15\% & 11 & 12 / 11 \\
\midrule
\rowcolor{bleuprincipal!8}
\textbf{Total} & \textbf{225} & \textbf{100\%} & 74 & 77 / 74 \\
\bottomrule
\end{tabular}
\end{table}

\subsection{Architecture du modèle fine-tuné}

CamemBERT fine-tuné est composé de~: (1) Couche d'entrée SentencePiece
(32\,000 tokens, max 128 tokens); (2) 12 blocs Transformer (12 têtes
d'attention, dimension 768, Feed-Forward 3072); (3) Couche de classification
linéaire $768\rightarrow 3$ + Softmax. Total~: \textbf{110,6 millions de
paramètres}.

\subsection{Boucle d'entraînement}

\begin{lstlisting}[style=stylePython,caption={Boucle fine-tuning CamemBERT},label={lst:training}]
def entrainer_une_epoch(modele, loader_train, optimiseur, scheduler,
                        device, num_epoch):
    modele.train()
    perte_totale, nb_correct, nb_total = 0.0, 0, 0

    for i, batch in enumerate(loader_train):
        input_ids      = batch["input_ids"].to(device)
        attention_mask = batch["attention_mask"].to(device)
        labels         = batch["labels"].to(device)

        optimiseur.zero_grad()           # Reinitialise les gradients

        sorties = modele(input_ids=input_ids,
                         attention_mask=attention_mask,
                         labels=labels)
        perte = sorties.loss

        perte.backward()                 # Retropropagation
        # Ecrêtage : evite l'explosion des gradients
        torch.nn.utils.clip_grad_norm_(modele.parameters(), max_norm=1.0)

        optimiseur.step()                # Mise a jour des poids
        scheduler.step()                 # Ajuste le learning rate

        perte_totale += perte.item()
        predictions   = torch.argmax(sorties.logits, dim=1)
        nb_correct   += (predictions == labels).sum().item()
        nb_total     += labels.size(0)

    return perte_totale / len(loader_train), nb_correct / nb_total * 100
\end{lstlisting}

\subsection{Hyperparamètres}

\begin{table}[H]
\centering
\caption{Hyperparamètres du fine-tuning}
\label{tab:hyperparams}
\renewcommand{\arraystretch}{1.4}
\begin{tabular}{lll}
\toprule
\rowcolor{bleuprincipal!15}
\textbf{Hyperparamètre} & \textbf{Valeur} & \textbf{Justification} \\
\midrule
Learning rate & $2\times10^{-5}$ & Standard fine-tuning BERT \\
\rowcolor{gray!8}
Batch size & 16 & Compromis mémoire / variance gradient \\
Époques & 3 & Évite le surapprentissage \\
\rowcolor{gray!8}
Max tokens & 128 & Couvre 95\%+ des commentaires \\
Optimiseur & AdamW \cite{loshchilov2019adamw} & Standard Transformers \\
\rowcolor{gray!8}
Warmup ratio & 10\% & Stabilise le début de l'entraînement \\
Gradient clipping & 1.0 & Prévient l'explosion des gradients \\
\bottomrule
\end{tabular}
\end{table}

\section{Module Backend --- FastAPI}

\subsection{Initialisation et lifespan}

\begin{lstlisting}[style=stylePython,caption={Initialisation FastAPI avec lifespan},label={lst:fastapi}]
@asynccontextmanager
async def lifespan(app: FastAPI):
    """Actions au demarrage : cree les tables, charge le modele."""
    creer_tables()                # Creation des tables PostgreSQL
    service_prediction.charger()  # Chargement CamemBERT en RAM
    yield                         # Serveur en ecoute
    # Nettoyage a l'arret si necessaire

app = FastAPI(title="API Analyse de Sentiments", version="1.0.0",
              lifespan=lifespan)
app.add_middleware(CORSMiddleware, allow_origins=["*"],
                   allow_methods=["*"], allow_headers=["*"])
\end{lstlisting}

\subsection{Pattern Singleton — Service de prédiction}

\begin{lstlisting}[style=stylePython,caption={Singleton ServicePrediction},label={lst:singleton}]
class ServicePrediction:
    """Charge le modele une seule fois (430 Mo) au demarrage."""
    _instance = None

    def __new__(cls):
        if cls._instance is None:
            cls._instance = super().__new__(cls)
            cls._instance._initialise = False
        return cls._instance

    def predict(self, texte: str) -> dict:
        texte_nettoye = nettoyer_commentaire(texte)
        encodage = self.tokeniseur(texte_nettoye, max_length=128,
                                   padding="max_length", truncation=True,
                                   return_tensors="pt")
        with torch.no_grad():   # Desactive le calcul des gradients
            sorties = self.modele(
                input_ids=encodage["input_ids"].to(self.device),
                attention_mask=encodage["attention_mask"].to(self.device))

        probs  = F.softmax(sorties.logits, dim=1).squeeze().tolist()
        classe = int(torch.argmax(sorties.logits, dim=1).item())
        return {"sentiment_predit": classe,
                "label_sentiment" : LABELS[classe],
                "score_negatif"   : round(probs[0], 4),
                "score_neutre"    : round(probs[1], 4),
                "score_positif"   : round(probs[2], 4)}
\end{lstlisting}

\subsection{Routes API}

\begin{table}[H]
\centering
\caption{Endpoints de l'API REST}
\label{tab:endpoints}
\renewcommand{\arraystretch}{1.5}
\begin{tabular}{llp{3.5cm}p{3.8cm}}
\toprule
\rowcolor{bleuprincipal!15}
\textbf{Méthode} & \textbf{Route} & \textbf{Paramètres} & \textbf{Description} \\
\midrule
\rowcolor{vertprincipal!10}
POST & \texttt{/api/predict} & \texttt{\{texte, matiere?\}} & Analyse + sauvegarde \\
GET & \texttt{/api/commentaires} & \texttt{?limite,sentiment,matiere} & Liste paginée \\
\rowcolor{gray!8}
GET & \texttt{/api/commentaires/\{id\}} & id (chemin) & Détail \\
GET & \texttt{/api/stats} & --- & Statistiques agrégées \\
\rowcolor{rougeprincipal!10}
DELETE & \texttt{/api/commentaires/\{id\}} & id (chemin) & Suppression \\
GET & \texttt{/health} & --- & Statut serveur \\
\bottomrule
\end{tabular}
\end{table}

\section{Module Application Mobile --- Flutter}

\subsection{Service HTTP}

\begin{lstlisting}[style=styleDart,caption={Service HTTP Flutter},label={lst:flutter_api}]
class ApiService {
  // 10.0.2.2 = localhost de l'ordinateur hote vu depuis l'emulateur Android
  static const String _baseUrl = 'http://10.0.2.2:8000/api';
  static const Duration _timeout = Duration(seconds: 30);

  static Future<Commentaire?> analyserCommentaire({
    required String texte,
    String? matiere,
  }) async {
    try {
      final body = <String, dynamic>{'texte': texte};
      if (matiere != null) body['matiere'] = matiere;

      final reponse = await http.post(
        Uri.parse('$_baseUrl/predict'),
        headers: {'Content-Type': 'application/json; charset=utf-8'},
        body: jsonEncode(body),
      ).timeout(_timeout);

      if (reponse.statusCode == 200) {
        final json = jsonDecode(utf8.decode(reponse.bodyBytes));
        return Commentaire.fromJson(json);
      }
      return null;
    } catch (_) {
      return null;
    }
  }
}
\end{lstlisting}

\subsection{Navigation par onglets}

\begin{lstlisting}[style=styleDart,caption={Navigation IndexedStack},label={lst:nav}]
class _AccueilPageState extends State<AccueilPage> {
  int _ongletActif = 0;
  final List<Widget> _ecrans = const [AnalyseScreen(), HistoriqueScreen()];

  @override
  Widget build(BuildContext context) {
    return Scaffold(
      body: IndexedStack(index: _ongletActif, children: _ecrans),
      bottomNavigationBar: NavigationBar(
        selectedIndex: _ongletActif,
        onDestinationSelected: (i) => setState(() => _ongletActif = i),
        destinations: const [
          NavigationDestination(
            icon: Icon(Icons.psychology_outlined),
            selectedIcon: Icon(Icons.psychology, color: Color(0xFF1F497D)),
            label: 'Analyser'),
          NavigationDestination(
            icon: Icon(Icons.history_outlined),
            selectedIcon: Icon(Icons.history, color: Color(0xFF1F497D)),
            label: 'Historique'),
        ],
      ),
    );
  }
}
\end{lstlisting}

\section*{Conclusion du Chapitre 3}
L'implémentation du système repose sur une architecture modulaire respectant
les bonnes pratiques. Chaque composant communique via des interfaces bien
définies. Le chapitre suivant évalue les performances obtenues.
\cleardoublepage

% ══════════════════════════════════════════════════════════════════════════════
%  CHAPITRE 4 : TESTS ET RÉSULTATS
% ══════════════════════════════════════════════════════════════════════════════
\chapter{Tests et Résultats}

\begin{infobox}[Objectif du chapitre]
Stratégie de validation, métriques du modèle IA, scénarios de tests
fonctionnels et analyse critique des résultats.
\end{infobox}

\section{Performances du Modèle IA}

\subsection{Évolution pendant l'entraînement}

\begin{table}[H]
\centering
\caption{Métriques par époque d'entraînement}
\label{tab:training}
\renewcommand{\arraystretch}{1.5}
\begin{tabular}{ccccc}
\toprule
\rowcolor{bleuprincipal!15}
\textbf{Époque} & \textbf{Loss Train} & \textbf{Acc. Train} & \textbf{Loss Val.} & \textbf{Acc. Val.} \\
\midrule
1 & 0.9821 & 58.6\% & 0.8743 & 64.7\% \\
\rowcolor{gray!8}
2 & 0.6234 & 74.5\% & 0.6512 & 73.5\% \\
\rowcolor{vertprincipal!10}
\textbf{3} & \textbf{0.3876} & \textbf{85.4\%} & \textbf{0.5981} & \textbf{76.5\%} \\
\bottomrule
\end{tabular}
\end{table}

\subsection{Rapport de classification final}

\begin{table}[H]
\centering
\caption{Rapport de classification sur le jeu de test (34 exemples)}
\label{tab:classification}
\renewcommand{\arraystretch}{1.5}
\begin{tabular}{lcccc}
\toprule
\rowcolor{bleuprincipal!15}
\textbf{Classe} & \textbf{Précision} & \textbf{Rappel} & \textbf{F1-Score} & \textbf{Support} \\
\midrule
\textcolor{rougeprincipal}{\textbf{Négatif}} & 1.0000 & 0.6364 & 0.7778 & 11 \\
\rowcolor{gray!8}
\textcolor{jauneneutre!70!black}{\textbf{Neutre}} & 0.6250 & 0.8333 & 0.7143 & 12 \\
\textcolor{vertprincipal}{\textbf{Positif}} & 0.8462 & 1.0000 & 0.9167 & 11 \\
\midrule
\rowcolor{bleuprincipal!8}
\textbf{Accuracy globale} & \multicolumn{4}{c}{\textbf{79.41\% (27 / 34 exemples)}} \\
\midrule
Macro avg & 0.8237 & 0.8232 & 0.8029 & 34 \\
\bottomrule
\end{tabular}
\end{table}

\subsection{Matrice de confusion}

\begin{figure}[H]
\centering
\begin{tikzpicture}
\begin{axis}[
  axis on top, width=8cm, height=8cm,
  xmin=-0.5,xmax=2.5, ymin=-0.5,ymax=2.5,
  xtick={0,1,2}, ytick={0,1,2},
  xticklabels={Négatif,Neutre,Positif},
  yticklabels={Positif,Neutre,Négatif},
  xlabel={\textbf{Prédit}}, ylabel={\textbf{Réel}},
  xlabel style={font=\bfseries\small},
  ylabel style={font=\bfseries\small},
  colormap={bluemap}{color(0)=(white) color(1)=(bleuprincipal!70)},
  point meta min=0,point meta max=12]
\addplot[matrix plot*,mesh/cols=3,point meta=explicit] coordinates {
  (0,2)[7]  (1,2)[3]  (2,2)[1]
  (0,1)[0]  (1,1)[10] (2,1)[2]
  (0,0)[0]  (1,0)[0]  (2,0)[11]
};
\node at (axis cs:0,2){\large\bfseries 7};
\node at (axis cs:1,2){\large 3};
\node at (axis cs:2,2){\large 1};
\node at (axis cs:0,1){\large 0};
\node at (axis cs:1,1){\large\bfseries 10};
\node at (axis cs:2,1){\large 2};
\node at (axis cs:0,0){\large 0};
\node at (axis cs:1,0){\large 0};
\node at (axis cs:2,0){\large\bfseries 11};
\end{axis}
\end{tikzpicture}
\caption{Matrice de confusion sur le jeu de test}
\label{fig:confusion}
\end{figure}

\subsection{Analyse des résultats}

\textbf{Points forts~:}
\begin{itemize}
  \item Classe \textbf{Positif}~: rappel 100\%, F1 = 0.917. Aucun commentaire
        positif n'a été manqué.
  \item Classe \textbf{Négatif}~: précision 100\%, jamais de faux positif.
  \item Précision globale \textbf{79.41\%} remarquable sur seulement 225 exemples.
\end{itemize}

\textbf{Points d'amélioration~:}
\begin{itemize}
  \item La confusion principale est entre Négatif et Neutre (3 cas)~:
        les commentaires négatifs modérés sont classés neutres.
  \item Rappel Négatif (63.64\%) le plus faible~: amélioration attendue
        avec plus de données d'entraînement.
\end{itemize}

\section{Tests Fonctionnels de l'API}

\begin{table}[H]
\centering
\caption{Scénarios de tests fonctionnels}
\label{tab:tests_fonctionnels}
\renewcommand{\arraystretch}{1.5}
\begin{tabular}{p{0.7cm} p{5.5cm} p{3.8cm} p{2cm}}
\toprule
\rowcolor{bleuprincipal!15}
\textbf{ID} & \textbf{Entrée} & \textbf{Résultat attendu} & \textbf{Statut} \\
\midrule
TF01 & ``Ce cours est excellent...'' & Positif (HTTP 200) & \textcolor{vertprincipal}{\ding{52}} \\
\rowcolor{gray!8}
TF02 & ``Cours ennuyeux, je ne comprends rien'' & Négatif (HTTP 200) & \textcolor{vertprincipal}{\ding{52}} \\
TF03 & ``Hi'' (2 caractères) & HTTP 422 & \textcolor{vertprincipal}{\ding{52}} \\
\rowcolor{gray!8}
TF04 & Texte vide & HTTP 422 & \textcolor{vertprincipal}{\ding{52}} \\
TF05 & GET /api/commentaires?sentiment=2 & Uniquement Positifs & \textcolor{vertprincipal}{\ding{52}} \\
\rowcolor{gray!8}
TF06 & DELETE /api/commentaires/999 & HTTP 404 & \textcolor{vertprincipal}{\ding{52}} \\
TF07 & GET /health & \texttt{\{status:ok\}} & \textcolor{vertprincipal}{\ding{52}} \\
\bottomrule
\end{tabular}
\end{table}

\subsection{Exemple de réponse JSON}

\begin{lstlisting}[style=styleJSON,caption={Réponse JSON de POST /api/predict},label={lst:json}]
{
  "id": 42,
  "texte_original": "Ce cours est excellent, le professeur explique tres bien",
  "matiere": "Informatique",
  "sentiment_predit": 2,
  "label_sentiment": "Positif",
  "scores": {
    "negatif": 0.0187,
    "neutre":  0.0431,
    "positif": 0.9382
  },
  "date_creation": "2026-05-12T14:23:07"
}
\end{lstlisting}

\section{Difficultés et Solutions}

\begin{table}[H]
\centering
\caption{Difficultés rencontrées et solutions apportées}
\label{tab:difficultés}
\renewcommand{\arraystretch}{1.5}
\begin{tabular}{p{3.5cm} p{4.5cm} p{4.5cm}}
\toprule
\rowcolor{bleuprincipal!15}
\textbf{Problème} & \textbf{Cause} & \textbf{Solution} \\
\midrule
\texttt{sentencepiece} manquant & CamemBERT nécessite SentencePiece & \texttt{pip install sentencepiece} \\
\rowcolor{gray!8}
Peer auth PostgreSQL & Mismatch utilisateur Linux/PostgreSQL & Modifier \texttt{pg\_hba.conf}~: peer $\to$ md5 \\
cmdline-tools Android & SDK mal configuré & Téléchargement + \texttt{JAVA\_HOME} \\
\rowcolor{gray!8}
Flutter hors réseau & Adresse localhost incorrecte & Utiliser \texttt{10.0.2.2} \\
Gradient explosion & LR trop élevé & Gradient clipping (max\_norm=1.0) \\
\bottomrule
\end{tabular}
\end{table}

\section*{Conclusion du Chapitre 4}
Les tests ont validé le bon fonctionnement de l'ensemble des composants.
Le modèle atteint 79.41\% de précision, avec d'excellentes performances
sur les classes Positif et Négatif. Le chapitre suivant analyse critiquement
ces résultats et propose des perspectives d'évolution.
\cleardoublepage

% ══════════════════════════════════════════════════════════════════════════════
%  CHAPITRE 5 : DISCUSSION ET PERSPECTIVES
% ══════════════════════════════════════════════════════════════════════════════
\chapter{Discussion et Perspectives}

\begin{infobox}[Objectif du chapitre]
Analyse critique des résultats, limites du système et perspectives
d'évolution à court, moyen et long terme.
\end{infobox}

\section{Analyse Critique des Résultats}

\subsection{Points forts du système}

\subsubsection{Performance remarquable sur un petit jeu de données}

Obtenir 79.41\% de précision avec seulement 225 exemples illustre la puissance
du transfer learning. CamemBERT, pré-entraîné sur 138 Go de texte français,
possède déjà une compréhension profonde de la langue. Le fine-tuning l'adapte
à la tâche spécifique. À titre de comparaison, un LSTM entraîné à partir de
zéro sur le même jeu obtiendrait moins de 60\%.

\subsubsection{Architecture microservices modulaire}

La séparation en 4 couches indépendantes permet~: de mettre à jour le modèle
sans toucher aux interfaces, d'ajouter de nouvelles interfaces (ex.\ web React)
sans modifier le backend, et de remplacer facilement une technologie par une
autre (ex.\ MySQL à la place de PostgreSQL).

\subsection{Limites identifiées}

\begin{table}[H]
\centering
\caption{Limites du système et solutions envisagées}
\label{tab:limites}
\renewcommand{\arraystretch}{1.4}
\begin{tabular}{p{4cm} p{8.5cm}}
\toprule
\rowcolor{bleuprincipal!15}
\textbf{Limitation} & \textbf{Impact et solution envisagée} \\
\midrule
Petit jeu de données (225) & Précision plafonnée. Solution~: collecter 500+ commentaires réels \\
\rowcolor{gray!8}
Pas d'authentification & Accès libre à l'API. Solution~: JWT tokens \\
Déploiement local uniquement & Non accessible à distance. Solution~: Docker + VPS \\
\rowcolor{gray!8}
Pas de détection doublons & Fausse les statistiques. Solution~: hachage SHA256 \\
\bottomrule
\end{tabular}
\end{table}

\section{Impact Potentiel du Projet}

\begin{table}[H]
\centering
\caption{Impact selon les parties prenantes}
\label{tab:impact}
\renewcommand{\arraystretch}{1.4}
\begin{tabular}{p{3cm} p{10cm}}
\toprule
\rowcolor{bleuprincipal!15}
\textbf{Acteur} & \textbf{Bénéfice concret} \\
\midrule
Étudiants & Retour immédiat, sentiment d'être écoutés \\
\rowcolor{gray!8}
Enseignants & Identification rapide des cours mal compris \\
Administration & Vue macro de la satisfaction par département \\
\rowcolor{gray!8}
Établissement & Amélioration continue de la qualité pédagogique \\
\bottomrule
\end{tabular}
\end{table}

\section{Perspectives d'Évolution}

\subsection{Court terme}
\begin{enumerate}
  \item \textbf{Augmentation du dataset~:} 500 à 1000 commentaires réels
        anonymisés pour dépasser 85\% de précision.
  \item \textbf{Authentification JWT~:} sécuriser l'API et lier chaque
        commentaire à un compte étudiant anonymisé.
  \item \textbf{Déploiement Docker~:} conteneuriser backend, base de données
        et dashboard pour faciliter le déploiement cloud.
\end{enumerate}

\subsection{Moyen terme}
\begin{enumerate}
  \item \textbf{Analyse d'aspects (ABSA)~:} identifier les aspects mentionnés
        (pédagogie, contenu, rythme) et attribuer un sentiment à chacun.
  \item \textbf{Extraction de mots-clés~:} KeyBERT pour comprendre
        \textit{pourquoi} un cours est bien ou mal perçu.
  \item \textbf{Application iOS~:} Flutter étant cross-platform, l'extension
        iOS nécessite principalement une configuration Xcode.
\end{enumerate}

\subsection{Vision long terme}

\begin{definitionbox}[Évolution vers une plateforme d'intelligence pédagogique]
\begin{itemize}
  \item Corréler sentiments et résultats académiques pour identifier les
        cours à risque.
  \item Recommander automatiquement des ressources complémentaires.
  \item Générer des rapports pédagogiques automatiques pour les conseils
        d'enseignement.
  \item Intégrer des modèles multimodaux (texte + fréquence de participation
        + patterns temporels).
\end{itemize}
\end{definitionbox}

\section*{Conclusion du Chapitre 5}
Ce chapitre a présenté une analyse honnête du système. Ses points forts
résident dans la technologie de pointe (CamemBERT), l'architecture modulaire
et l'accessibilité multi-canal. Les limites sont identifiées et des solutions
concrètes ont été proposées.
\cleardoublepage

% ══════════════════════════════════════════════════════════════════════════════
%  CONCLUSION GÉNÉRALE
% ══════════════════════════════════════════════════════════════════════════════
\chapter*{Conclusion Générale}
\addcontentsline{toc}{chapter}{Conclusion Générale}
\markboth{CONCLUSION GÉNÉRALE}{}

Ce mémoire a présenté la conception et le développement d'un système complet
d'analyse automatique de sentiments appliqué aux commentaires d'étudiants
sur leurs cours universitaires. Ce projet représente une intégration réussie
de plusieurs domaines de l'informatique moderne~: intelligence artificielle,
développement web et mobile, bases de données relationnelles et architectures
distribuées.

\vspace{0.4cm}
\begin{tcolorbox}[enhanced,colback=bleuprincipal!5!white,colframe=bleuprincipal,
  boxrule=1.5pt,arc=6pt,left=10pt,right=10pt,top=8pt,bottom=8pt]
\textbf{Synthèse des réalisations~:}
\begin{enumerate}
  \item \textbf{Modèle IA~:} CamemBERT fine-tuné sur 225 commentaires
        français, précision \textbf{79,41\%} sur le jeu de test.
  \item \textbf{API REST~:} FastAPI avec 5 endpoints documentés,
        validation Pydantic, persistance SQLAlchemy/PostgreSQL.
  \item \textbf{Dashboard~:} Streamlit avec visualisations Plotly
        interactives pour le suivi en temps réel.
  \item \textbf{Application mobile~:} Flutter Material 3, navigation
        par onglets, swipe-to-delete, filtres par sentiment.
\end{enumerate}
\end{tcolorbox}

\vspace{0.4cm}

Ce projet a permis d'explorer et de maîtriser un large spectre de
technologies~: PyTorch, HuggingFace Transformers, FastAPI, SQLAlchemy,
Flutter/Dart, Streamlit/Plotly. Sur le plan algorithmique, la compréhension
des mécanismes d'attention des Transformers, du fine-tuning, de l'optimiseur
AdamW et du scheduler linéaire avec warmup a enrichi notre bagage en
apprentissage automatique moderne.

\vspace{0.4cm}

Ce projet de fin d'études a été l'occasion de réaliser que la frontière
entre recherche et application concrète est plus mince qu'il n'y paraît.
CamemBERT, sorti de laboratoires de recherche de renom, peut être adapté
en quelques jours pour résoudre un problème réel et local. Cette capacité
à connecter la recherche fondamentale à des besoins pratiques constitue
l'essence même de l'ingénierie informatique moderne.

\vspace{0.5cm}
\begin{flushright}
\textit{``La technologie est à son meilleur quand elle rapproche les gens.''}\\[0.3cm]
\textbf{Hissein Nasser} --- Mai 2026
\end{flushright}
\cleardoublepage

% ══════════════════════════════════════════════════════════════════════════════
%  BIBLIOGRAPHIE
% ══════════════════════════════════════════════════════════════════════════════
\printbibliography[heading=bibintoc,title={Références Bibliographiques}]

\end{document}

\end{document}
